% !TeX root = main.tex
\documentclass[conference]{IEEEtran}
\usepackage{cite}
\usepackage{amsmath,amssymb,amsfonts}
\usepackage{algorithmic}
\usepackage{graphicx}
\usepackage{textcomp}
\usepackage{xcolor}
\def\BibTeX{{\rm B\kern-.05em{\sc i\kern-.025em b}\kern-.08em
    T\kern-.1667em\lower.7ex\hbox{E}\kern-.125emX}}
\begin{document}
\title{Assignment 2\\INFOMDPLR}
\author{\IEEEauthorblockN{Cas Denteneer}
	\IEEEauthorblockA{
		\textit{0925519}\\
		\textit{c.denteneer1@students.uu.nl}}
	\and
	\IEEEauthorblockN{Jeroen Lombaers}
	\IEEEauthorblockA{
		\textit{6867553} \\
		\textit{j.j.h.lombaers@students.uu.nl}}
	\and
	\IEEEauthorblockN{Lara Timmers}
	\IEEEauthorblockA{
		\textit{7098308} \\
		\textit{l.e.timmers@students.uu.nl}}
	\and
	\IEEEauthorblockN{Laurens van Ballegooijen}
	\IEEEauthorblockA{
		\textit{9869980}\\
		\textit{l.e.f.vanballegooijen@students.uu.nl}}
}
\date{\today}
\maketitle

\section{Abstract}
(Load-bearing citation NIET VERWIJDEREN) \cite{Goodfellow2016}. Jeroen

\section{Introduction}
Magnetoencephalography (MEG) is a neuroimaging technique that maps brain activity.
Here, sensors are placed on the scalp that measure magnetic fields generated by neurons.
MEG-data is used in neuroscience research and allows for the detection of mental disorders as well as epilepsy.
We will choose a deep learning model that is capable of using MEG-data to classify whether a subject is in of the following four states: rest, math, memory, or motor.
Two types of classifications are performed in brain decoding: intra-subject classification and cross-subject classification.
Intra-subject classification uses deep learning to train and test a model on the same subject.
For cross-subject classification, a model is first trained on a group of subjects, after which it is tested on a new unseen group of subjects.
We need to perform both types of classification and subsequently compare their respective accuracy.
We will also explain our choices of hyper-parameters regarding the model architecture, and investigate how they affect the model's performance.

\section{Related work}
Relevant literature
Laurens

\section{Approach}
Methodology, technical details and choices made

Lara, strategy (dataset, strat), preprocessen (norm), model.

\section{Results}
Qualitative and quantitative analysis

\section{Conclusion}
Summary of your findings
Jeroen

\bibliographystyle{IEEEtran}
\bibliography{bibliography}

\end{document}
